%%%%%%%%%%%%%%%%%%%%%%%%%%%%%%%%%%%%%%%%%%%%%%%%%%%%%%%%%%%%%%%%%%%%%%%%%%%%%%%%%%%%%%%%%%%%%%%%%%% 
% KAPITOLA 6
%%%%%%%%%%%%%%%%%%%%%%%%%%%%%%%%%%%%%%%%%%%%%%%%%%%%%%%%%%%%%%%%%%%%%%%%%%%%%%%%%%%%%%%%%%%%%%%%%%
\chapter{Implementácia}
\label{ch:kap6}
Táto kapitola sa zameriava na implementáciu predstaveného návrhu v Kapitole~\ref{ch:kap5}. Nezaoberá sa  opisom zdrojového kódu, ale implementačnými detailmi, problémami ktoré sa pri ich implementácii objavili, a spôsob ich riešenia. Opisuje použité technológie, ako Python, SQLAlchemy, Ollama, atď (Sekcia~\ref{sub:pouzite_technologie}).  Ďalej vysvetľuje akým spôsobom je získavaná štruktúra databázy a vzorové dáta z jej tabuliek (Sekcia~\ref{sub:extrakcia_a_vzorky}). Vysvetľuje rozdelenie celého procesu generovania do štyroch samostatných fáz (Sekcia~\ref{sub:fazy_generovania_receptu}) a špecializáciu lokálnych modelov pomocou textových súborov Modelfile (Sekcia~\ref{sub:implementacia_modelfiles}) a mechanizmus validácie a opráv výstupov (Sekcia~\ref{sub:validacia_oprava}). Záver kapitoly opisuje výstup nástroja (Sekcia~\ref{sub:vystupy}).

\section{Použité technológie}
\label{sub:pouzite_technologie}

\begin{itemize}
    \item \textbf{Python 3.12} implementačný jazyk nástroja
    \item \textbf{SQLAlchemy} extrakcia schémy databázy nezávisle od DBMS (Sekcia~\ref{sub:extrakcia_a_vzorky})
    \item \textbf{PyYAML} spracovanie receptov vo formáte YAML, vrátane vlastného loaderu \newline \texttt{UniqueKeyLoader} na detekciu duplicitných kľúčov v schéme
    \item \textbf{Ollama} lokálne spúšťanie LLM a ich prispôsobenie cez Modelfiles (Sekcia~\ref{sub:implementacia_modelfiles})
    \item \textbf{qwen2.5:3b/7b~\cite{Qwen2025}, llama3.1:8b~\cite{Llama2024}, mistral:7b~\cite{Mistral7b}} použité jazykové modely (Sekcia~\ref{sub:testovacie_scenare})
    \item \textbf{Dataster/dbgenx} cieľový generátor testovacích dát, recepty sú overené priamym volaním knižnice dbgenx, nad ktorou je Dataster postavený
    \item \textbf{MySQL} referenčná DBMS
    \item \textbf{PostgreSQL} DBMS pre demonštráciu prenosnosti
    \item \textbf{Docker} izolované testovacie databázy
    
\end{itemize}

\section{Extrakcia a vzorky}
\label{sub:extrakcia_a_vzorky}

V tejto časti pipeline stojí kód medzi databázou a prvou fázou LLM, zbiera všetky potrebné informácie z databázy pre generovanie a skladá ich do serializovanej štruktúry, ktorá je vstupom do \texttt{types} fázy (\ref{sub_sub:types}). Táto implementácia  sa zaoberá mechanickým získavaním informácií.


Štruktúra schémy sa získava prostredníctvom SQLAlchemy Inspector v module \newline schema\_extractor.py. Umožňuje to získavať informácie naprieč rôznymi databázami bez nutnosti upravovania SQL dopytov aby zodpovedali konkrétnej {DBMS}. V počiatočných fázach vývoja bol najskôr implementovaný len MySQL konektor. Tento prístup sa ukázal ako limitujúci, keďže implementácia nepodporovala ďalšie {DBMS} a tým pádom bol použitý SQLAlchemy Inspector.


Ďalšou dôležitou časťou sú samotné dáta posielané do generujúcich fáz. Bez nich nie je možné vytvoriť sémanticky zmysluplný recept. Stratégia zberu vzoriek sa delí podľa charakteru stĺpca:

\begin{itemize}
    \item \textbf{Stĺpce primárneho kľúča (PK):} pri primárnych kľúčoch sa nezbierajú vzorky rovnako ako pri ostatných stĺpcoch keďže primárny kľúč je ľahko odvoditeľný. Namiesto toho sa zbierajú minimálne a maximálne hodnoty, aby bol určený presný rozsah primárneho kľúča, ktorý sa následne obalí do konštrukcie \texttt{seq}, ako vidno v ukážke. 
    \item \textbf{Ostatné stĺpce (Non–PK):} pri ostatných stĺpcoch sa používa SQL dopyt využívajúci operáciu UNION, kde sa vyberajú unikátne hodnoty vzostupne a zostupne. Dôvodom je aby sa v prípade rôznych intervalov zobrazovali skutočné hraničné hodnoty. Toto nám dáva presný obraz o dátach v tabuľkách, ktoré sú kľúčové pre generovanie types (hlavne intervalov) a obmedzení v constraints.
    \item \textbf{Filtrovanie hodnôt NULL:} všetky SQL dopyty na vzorky obsahujú podmienku \texttt{IS NOT NULL}, ktorá odfiltruje prázdne hodnoty, pretože cieľom je získať dáta ktoré nám reálne priblížia sémantiku hodnôt stĺpca.
    \item \textbf{Limitácia vzoriek:} množstvo vzoriek je limitované funkciou \texttt{calculate\_limit}, ktorá na základe celkového počtu riadkov v tabuľke vypočíta ideálny počet vzoriek, okolo 25\% celkových dát, na rozsah okolo 5-20 vzoriek. Do promptu sa následne odovzdáva najviac 10 vybraných hodnôt spolu s hodnotou \texttt{distinct\_count\_in\_sample}, ktorá udáva počet unikátnych hodnôt v tejto odobranej vzorke. Tento limit udržiava prompt v rozumnej veľkosti, ale stále poskytuje dosť vzoriek pre zachytenie rozsahu a charakteru hodnôt.
\end{itemize}

Po získaní štruktúry a dát je potrebné poskladať informácie do jedného spoločného formátu JSON, niektoré relačné databázy ale obsahujú typy, ktoré nie je možné transformovať do JSON formátu a preto sa používa pomocná funkcia json\_decimal(), ktorá transformuje typy ako \texttt{DECIMAL} na float, \texttt{timedelta} na \texttt{HH:MM:SS} a hodnoty \texttt{datetime/date} na reťazec oddelený medzerou namiesto štandardného ISO znaku \texttt{T}.
\begin{figure}[H]
    \centering
    \begin{lstlisting}[caption={Príklad serializovaného vstupu}, label={lst:json}, frame=single, linewidth=0.9\textwidth, xleftmargin=0.09\textwidth]
    {
      "users": {
        "id": {
          "is_pk": true,
          "db_type": "int",
          "data_type": "numeric",
          "seq": "seq(1, 10, 1)"
        },
        "name": {
          "is_pk": false,
          "db_type": "varchar(100)",
          "distinct_count_in_sample": 5,
          "samples": [
            "Aneta Valekova",
            "Eva Vesela",
            "Frantisek Kucera",
            "Pavel Cerny",
            "Stepan Horak"
          ]
        },
        "email": {
          "is_pk": false,
          "db_type": "varchar(150)",
          "distinct_count_in_sample": 5,
          "samples": [
            "aneta@test.com",
            "eva@gmail.com",
            "frantisek@test.com",
            "peter@test.com",
            "stepan@bing.com"
          ]
        }
      }
    }
\end{lstlisting}
\end{figure}

Celý tento proces je vykonaný bez využitia jazykového modelu, tým sa predišlo možným halucináciám a nepredvídateľnému správaniu, ktoré by mohlo nastať, ak by bol jazykový model implementovaný. Po vytvorení serializovaného vstupu pre types fázu úloha tejto časti tým pádom končí a výsledok je odovzdaný ako vstup pre fázy generovania.

\section{Fázy generovania receptu}
\label{sub:fazy_generovania_receptu}

Proces generovania receptu sa delí do štyroch fáz: \texttt{types}, \texttt{schema}, \texttt{generate} a \texttt{constraints}. Všetky fázy zdieľajú spoločnú funkciu pre komunikáciu s rozhraním Ollama klienta.

Pre volanie klienta boli definované aj nasledujúce parametre:

\begin{itemize}
    \item \textbf{Kontextové okno (\texttt{num\_ctx)}} - bolo nastavené na 16384 tokenov. Veľké kontextové okno bolo potrebné pre rozsiahle serializované schémy so vzorkami dát pri databázach, ktoré majú desiatky tabuliek.
    \item \textbf{Maximálny počet vygenerovaných tokenov (\texttt{num\_predict)}}: bol nastavený na 8192 tokenov, aby sa predišlo predčasnému orezaniu sekcie
    \item \textbf{Časový limit (\texttt{timeout)}}: nastavený na 600 sekúnd (10 minút) na jedno volanie, tým sa zabráni zaseknutiu procesu.  
\end{itemize}

V tomto projekte je limitácia 8 GB VRAM. Fázy s interpretačne náročnejšou úlohou, ako constraints a types, obsluhuje väčší model (qwen2.5:7b), ľahšie fázy variant qwen2.5:3b. Zmiešaný prístup modelov, qwen2.5:7b a qwen2.5:3b, umožňuje využitie väčšieho modelu tam kde úloha vyžaduje interpretáciu a menší tam kde stačí iba transformácia. Porovnávaním jednotlivých rodín a veľkosti modelov pomocou testov sa zaoberá Kapitola~\ref{ch:kap8}. Na záver všetky fázy sú poskladané do jedného receptu a prechádzajú postupnosťou opravy a validácie.

\subsection{Types}
\label{sub_sub:types}
Vstupom fázy types je serializovaná schéma so vzorkami dát, ich extrakciu opisuje Sekcia~\ref{sub:extrakcia_a_vzorky}. Model každému stĺpcu priradí dátový typ v súlade s nástrojom Dataster, napríklad \texttt{seq(min,max,step)} pre primárne kľúče, \texttt{interval/finterval/dinterval} s hranicami zo vzoriek, \texttt{str\_len} alebo \texttt{regex} pre reťazce a referencie na datasety. Výstup fázy types priamo ovplyvňuje tvorbu a kvalitu ostatných sekcií. Je priamou súčasťou promptov všetkých troch nasledujúcich fáz.  \texttt{Schema}, \texttt{generate} aj \texttt{constraints} ho dostávajú ako vstup. Akákoľvek chyba sa preto rozšíri priamo do ostatných fáz. Samostatná fáza umožňuje vyhraničenie úlohy pre malý model, špecializovanými few-shot príkladmi a možnosťou validovať a opravovať výstupy jednotlivých fáz oddelene. Výstup sekcie je surový text sekcie types, ktorý sa validuje (Sekcia~\ref{sub:validacia_oprava}). Jedným konkrétnym príkladom opravy je nesprávne priradenie dátového typu modelom pri stĺpcoch, ktoré obsahovali DECIMAL hodnoty. Príčinou tohto problému bolo, že Python DECIMAL sa vložil do serializovanej schémy ako reťazec a tým pádom model interpretoval tieto hodnoty ako reťazcový typ namiesto číselného. Riešením je konverzia DECIMAL na float pri serializácii vstupu.

\subsection{Schema}
\label{sub_sub:schema}
Vstupom fázy \texttt{schema} je prompt, ktorý je zložený z výstupu sekcie \texttt{types}, kritických pravidiel a validných príkladov. Kritické pravidlá stanovujú to, že každá tabuľka a jej stĺpce sa musia objaviť v schéme, ďalej sú tam pravidlá formátovania a zákaz halucinovania hodnôt, ktoré neexistujú. Bez týchto pravidiel model vynecháva stĺpce alebo vytvára také, ktoré sa v tabuľke nenachádzajú. Úlohou tejto sekcie je mapovanie stĺpcov na atribúty typov z predchádzajúcej fázy (\texttt{types}). Zároveň slúži na vyjadrenie cudzích kľúčov pomocou štruktúry \texttt{FROM}.  Konkrétny príklad opravy je zlyhanie LLM pri návrate nesprávneho názvu stĺpcov, samotný názov sa síce v tabuľke objavil, lenže pri určitej DBMS (napríklad PostgreSQL), ktorá má citlivosť na malé a veľké písmená, došlo k chybe. 

\subsection{Generate}
\label{sub_sub:generate}
Vstupom fázy \texttt{generate} je prompt poskladaný z \texttt{types} sekcie vygenerovanej z predchádzajúcej fázy, kritických pravidiel a platných príkladov. Kritické pravidlá upozorňujú na správny formát sekcie \texttt{generate} a zároveň určujú, že pre každú tabuľku musí byť vygenerovaný počet riadkov. Posledné pravidlo je, že počty riadkov musia odzrkadľovať hodnoty vo vnútri konštrukcie seq, inak sa tam dosadí generické číslo. Bez týchto pravidiel model vytvára hodnoty, ktoré môžu spôsobiť zlyhanie pri generovaní testovacích dát. 

\subsection{Constraints}
\label{sub_sub:constraints}
Úlohou tejto fázy je interpretovať význam hodnôt v \texttt{types} sekcii a vytvoriť logické obmedzenia. Vstupom fázy je prompt zložený z \texttt{types} výstupu, prísnych pravidiel a platného príkladu. 

Pravidlá obmedzení hovoria o tom, že pri generovaní obmedzení sa majú používať len platné názvy tabuliek a ich stĺpcov, ktoré sú uvedené v predtým vygenerovanej časti \texttt{types}. Každé obmedzenie sa zapisuje na nový riadok so špecifickým formátom, ďalej určuje nepovolené operátory a zakazuje generovanie obmedzení s referenciami na samého seba. Táto fáza je spoločne so sekciou \texttt{types} jediná interpretačná úloha. Jedným z problémov pri generovaní \texttt{constraints} bola referencia na samého seba \newline (\textit{users.age = users.age}). Toto obmedzenie nepridáva do výsledného receptu žiadnu hodnotu a preto sú tieto obmedzenia vynechané.



\section{Modelfiles}
 \label{sub:implementacia_modelfiles}
Modelfile je textový súbor, ktorý slúži ako podklad pre vytváranie a prispôsobenie LLM pre mechanizmus Ollama. Každá fáza má vlastný špecializovaný Modelfile so systémovým promptom a few-shot príkladmi. Zatiaľ čo Modelfile má v sebe pravidlá správania a few-shot príklady, prompty zo súboru \textbf{prompts.py} mu dodávajú serializované schémy konkrétnej databázy.


Každý Modelfile má podobnú štruktúru, teda tri základné bloky:
\begin{itemize}
    \item \textbf{FROM:} týmto sa určuje tak vybraná rodina modelov, ako aj veľkosť a kapacita modelu. 
    \item \textbf{SYSTEM:} obsahuje všeobecné inštrukcie, pravidlá a logické obmedzenia, ktoré musí model dodržiavať
    \item \textbf{Few-shot MESSAGE user/assistant:} slúži ako podklad konverzácií medzi používateľom a asistentom. Model ich používa ako vzor správania.
\end{itemize}

Všeobecné inštrukcie vo vnútri \textbf{SYSTEM} promptu vznikali postupne počas testovania programu. Každé pravidlo predstavuje nejakú chybu, alebo moment keď sa model odchýlil od inštrukcií. Modelfile preto obsahuje explicitné zákazy a pravidlá.

\begin{figure}[H]
    \centering
    \begin{lstlisting}[caption={Príklad pravidla}, label={lst:rule_example}, frame=single, linewidth=0.9\textwidth, xleftmargin=0.05\textwidth]
ALWAYS copy column names EXACTLY as they are spelled in the 
input JSON database samples. Do NOT fix typos in column names
(for example., if input says 'recieved_date', output MUST be 
'recieved_date').
\end{lstlisting}

\end{figure}

Few-shot je história konverzácie vložená do každého Modelfile. Príklady boli vybrané tak, aby pokryli rôzne problémové scenáre, ktoré boli objavené počas testovania. Pokrývajú prípady ako vytvorenie \textbf{BLOCKS} formátu, generovanie poskladaných hodnôt, atď.

\begin{figure}[H]
    \centering
    \begin{lstlisting}[caption={Príklad BLOCKS few-shot správy}, label={lst:fewshot_example}, frame=single, linewidth=0.95\textwidth, xleftmargin=0.05\textwidth]
MESSAGE user """Database samples:
{"Currencies": 
    {"currency_code": 
        {"is_pk": true, 
         "db_type": "varchar(3)", 
         "data_type": "non-numeric", 
         "samples": ["CZK", "EUR", "GBP", "USD"]}, 
     "currency_name": 
        {"is_pk": false, 
         "db_type": "varchar(50)", 
         "distinct_count_in_sample": 2, 
         "samples": ["Czech Koruna", "Euro"]}}}

Return ONLY the types section starting with 'types:', nothing else. 
Do not wrap code in markdown block format."""

MESSAGE assistant """
    types:
      Currencies:
        currency_code:
          - BLOCKS
          - CZK
          - EUR
          - GBP
          - USD
        currency_name: str_len([cCs] 4,50)
"""

\end{lstlisting}

\end{figure}

Aj napriek všetkým pravidlám sa môže stať, že ich model poruší. Práve preto existuje validačná a opravná vrstva.

\section{Validácia a oprava}
\label{sub:validacia_oprava}

Výstup jazykového modelu je nedeterministický a ani prísne prompty nedokážu zaručiť validný výstup. Preto po poskladaní celého receptu nasleduje automatická oprava a validácia. Namiesto opakovaného volania modelu sa opravy vykonávajú v kóde. Dôvodom je, že oprava v kóde je reprodukovateľná a výhodnejšia ako opakované volanie modelu. Bola vytvorená postupnosť opráv, ktorá vygenerovaný recept postupne opravuje. Oprava sa spúšťa až nad poskladaným receptom.

Priebeh validácie zloženého receptu znázorňuje obrázok~\ref{fig:opravavalidacia}. Ak je YAML vygenerovaný bez chýb, recept sa úspešne uloží. V prípade, že sa v recepte objaví chyba, tak sa zaznamená a recept sa taktiež uloží. V prípade, že sa v recepte nachádza chyba, ktorú nie je možné opraviť, dochádza ku zlyhaniu (Kapitola~\ref{ch:kap7}).

\begin{figure}[H]
    \centering
    \includegraphics[width=1\linewidth]{obrazky-figures/opravavalidacia.pdf}
    \caption{Dvojvrstvový mechanizmus opráv a záverečná validácia}
    \label{fig:opravavalidacia}
\end{figure}


Dôvod týchto opráv je jednoduchý.  V prípade, že by sa chyba objavila v sekcii \texttt{types} a zostala neopravená, bola by okamžite poslaná do promptov ostatných sekcií, ktoré ju nasledujú.
 
Hodnotové opravy predpokladajú korektnú štruktúru dokumentu, práve preto sú opravy rozdelené na dve vrstvy:

\begin{itemize}
    \item \textbf{štrukturálna oprava} (napríklad list-of-dict)
    \item \textbf{hodnotové opravy} (zlé názvy, typy)
\end{itemize}

Toto rozdelenie je z dôvodu, že ak dôjde k deformácii formátu YAML, hodnotová kontrola sa buď vôbec nespustí alebo sa vykonajú opravy na zlom mieste.


V prípade, že sa ale vyskytne chyba, ktorú kód nevie opraviť, napríklad rozpadnutá štruktúra YAML, zapíše sa to ako chybový stav do data class, kritické chyby sa uložia do CSV správy a generovanie pre danú databázu sa ukončí. Pôvodný poškodený súbor sa zapíše do results/original\_yaml.


Táto práca sa zameriavala na to, aby bola generácia YAML nielen správna, ale aj čo najviac použiteľná v praxi. Preto bolo implementované takzvané zrkadlenie opráv, ktoré zabezpečuje prepojenie medzi validačnou a opravnou vrstvou. Ukončenie procesu na validácii by znamenalo zbytočné zahadzovanie vygenerovaných receptov, ktoré boli v jadre správne, lenže obsahovali drobnú chybu. Keďže kód už mal mechanizmus na detekciu týchto chýb, ich automatická oprava bola logickým a priamym pokračovaním existujúcej validácie.

Čo sa týka opráv, sú tu aj určité limitácie. V prípade, že LLM vygeneruje neznámy alebo nezmyselný dátový typ, tak sa kód pokúsi o jeho záchranu pomocou funkcie \newline \textit{guess\_type\_from\_db}. Avšak, ak to zlyhá, kód ako poslednú možnosť zmení typ daného stĺpca na univerzálny typ \textbf{string}. Aj keď táto oprava zachráni recept pred zlyhaním, samotná kvalita týmto klesá. Vygenerované hodnoty strácajú sémantický význam a vzťah k reálnym dátam v databáze. Každý zásah je zaznamenaný v logu a štatistikách o opravách.

\section{Výstup}
\label{sub:vystupy}
Po prejdení všetkých fáz sa poskladá konečný recept a dopĺňa sa \texttt{include} sekcia, ktorá sa generuje deterministicky zo zoznamu použitých datasetov. Finálny recept sa ukladá do \textbf{results/dataster}, surový recept pred opravami sa ukladá do \textbf{results/original\_yaml/} a CSV so štatistikami chýb a opráv sa ukladá do \textbf{results/csv/}.


